\documentclass[11pt]{article}
\usepackage{amsmath,amssymb,amsthm}
\usepackage[margin=1in]{geometry}
\title{Problem 705: Total Inversion Count of Digit Strings}
\author{Project Euler Solutions}
\date{}

\begin{document}
\maketitle

\section{Problem Statement}
For a digit string $d_1 d_2 \cdots d_k$, the inversion count is the number of pairs $(i,j)$ with $i < j$ and $d_i > d_j$. Let $T(n)$ be the total inversion count over all $n$-digit strings (digits 0--9). Find $T(10^{16}) \bmod 10^9 + 7$.

\section{Mathematical Analysis}
Each $n$-digit string uses digits $\{0,1,\ldots,9\}$. For a fixed pair of positions $(i,j)$ with $i<j$, the expected number of inversions is the probability that a random digit is greater than another random digit.

For two independent uniform digits from $\{0,\ldots,9\}$, the probability of inversion is $\frac{45}{100}$ (there are $\binom{10}{2}=45$ ordered pairs with $d_i > d_j$ out of $100$ total).

\section{Derivation}
The number of position pairs is $\binom{n}{2}$. Each pair contributes $45$ inversions on average across all $10^n$ strings (since there are $45 \cdot 10^{n-2}$ strings where that specific pair is inverted). The total inversion count is:

$$T(n) = \binom{n}{2} \cdot 45 \cdot 10^{n-2}$$

$$= \frac{n(n-1)}{2} \cdot 45 \cdot 10^{n-2} = \frac{45 \cdot n(n-1)}{2} \cdot 10^{n-2}$$

For $n = 10^{16}$, compute modulo $10^9+7$.

\section{Proof of Correctness}
By linearity of summation, the total inversions across all strings equals the number of (string, position-pair) combinations where an inversion occurs. Each position pair $(i,j)$ sees exactly $45 \cdot 10^{n-2}$ inversions (choose the ordered pair of digits in 45 ways, and the remaining $n-2$ digits in $10^{n-2}$ ways).

\section{Complexity Analysis}
$O(\log n)$ using modular exponentiation.

\section{Answer}

\[
\boxed{480440153}
\]

\end{document}
